\documentclass[11pt]{article}

\usepackage[T1]{fontenc}
\usepackage[utf8]{inputenc}
\usepackage{lmodern}
\usepackage{microtype}
\usepackage[margin=1in]{geometry}
\usepackage{amsmath,amssymb,amsthm,mathtools,bm}
\usepackage{mathrsfs}
\usepackage{booktabs,array,longtable}
\usepackage{enumitem}
\usepackage{xcolor}
\usepackage[most]{tcolorbox}
\usepackage{fancyhdr}
\usepackage{hyperref}
\usepackage[nameinlink,capitalise,noabbrev]{cleveref}

\definecolor{deepblue}{RGB}{25,61,105}
\definecolor{deepred}{RGB}{132,32,41}
\definecolor{deepgreen}{RGB}{28,92,68}
\definecolor{softblue}{RGB}{239,246,253}
\definecolor{softred}{RGB}{253,241,242}
\definecolor{softgreen}{RGB}{239,249,245}
\definecolor{softgray}{RGB}{246,247,249}

\hypersetup{
  colorlinks=true,
  linkcolor=deepblue,
  citecolor=deepgreen,
  urlcolor=deepblue,
  pdftitle={Alaoglu--Erdos at Depth Two},
  pdfauthor={Independent continuation report},
  pdfsubject={Alaoglu--Erdos conjecture, four exponentials, logarithms of algebraic numbers}
}

\pagestyle{fancy}
\fancyhf{}
\fancyhead[L]{\small Alaoglu--Erd\H{o}s at depth two}
\fancyhead[R]{\small Independent continuation report}
\fancyfoot[C]{\thepage}
\setlength{\headheight}{14pt}

\setcounter{tocdepth}{2}
\setcounter{secnumdepth}{3}
\setlist[itemize]{leftmargin=2em,itemsep=2pt,topsep=4pt}
\setlist[enumerate]{leftmargin=2.4em,itemsep=2pt,topsep=4pt}

\newtheorem{theorem}{Theorem}[section]
\newtheorem{proposition}[theorem]{Proposition}
\newtheorem{lemma}[theorem]{Lemma}
\newtheorem{corollary}[theorem]{Corollary}
\newtheorem{conjecture}[theorem]{Conjecture}
\newtheorem{question}[theorem]{Question}
\theoremstyle{definition}
\newtheorem{definition}[theorem]{Definition}
\newtheorem{example}[theorem]{Example}
\theoremstyle{remark}
\newtheorem{remark}[theorem]{Remark}
\newtheorem{warning}[theorem]{Warning}

\newtcolorbox{auditbox}[1][]{
  enhanced,breakable,colback=softred,colframe=deepred,
  fonttitle=\bfseries,title=Audit conclusion,#1
}
\newtcolorbox{resultbox}[1][]{
  enhanced,breakable,colback=softgreen,colframe=deepgreen,
  fonttitle=\bfseries,title=New unconditional result,#1
}
\newtcolorbox{frontierbox}[1][]{
  enhanced,breakable,colback=softblue,colframe=deepblue,
  fonttitle=\bfseries,title=Exact frontier,#1
}
\newtcolorbox{plainbox}[1][]{
  enhanced,breakable,colback=softgray,colframe=black!45,
  fonttitle=\bfseries,#1
}

\newcommand{\Q}{\mathbb{Q}}
\newcommand{\Z}{\mathbb{Z}}
\newcommand{\N}{\mathbb{N}}
\newcommand{\R}{\mathbb{R}}
\newcommand{\C}{\mathbb{C}}
\newcommand{\Qbar}{\overline{\mathbb{Q}}}
\newcommand{\Apos}{\overline{\mathbb{Q}}_{>0}}
\newcommand{\Nzero}{\mathbb{N}_{0}}
\newcommand{\Npos}{\mathbb{N}_{>0}}
\newcommand{\Span}{\operatorname{span}}
\newcommand{\rank}{\operatorname{rank}}
\newcommand{\Sym}{\operatorname{Sym}}
\newcommand{\supp}{\operatorname{supp}}
\newcommand{\ord}{\operatorname{ord}}
\newcommand{\diag}{\operatorname{diag}}
\newcommand{\Gm}{\mathbb{G}_{m}}
\newcommand{\PGL}{\operatorname{PGL}}
\newcommand{\FE}{\mathrm{FE}_4}
\newcommand{\MCC}{\mathrm{MCC}_2}
\newcommand{\SRC}{\mathrm{SRC}}
\newcommand{\AEconj}{\mathrm{AE}_{2,3}}
\newcommand{\ellmap}{\boldsymbol{\ell}}
\newcommand{\vct}[1]{\boldsymbol{#1}}
\newcommand{\ip}[2]{\left\langle #1,#2\right\rangle}

\title{\bfseries Alaoglu--Erd\H{o}s at Depth Two\\[3pt]
\Large An Independent Proof Audit, Quadratic-Rank Reduction,\\
and the Exact Remaining Obstruction}
\author{Independent continuation report}
\date{21 August 2026}

\begin{document}
\maketitle

\begin{abstract}
We investigate the assertion
\[
   2^x,3^x\in\Z \quad\Longrightarrow\quad x\in\Z
\]
for real $x$.  The supplied unified report already develops a very large collection of
semigroup, valuation, determinant, approximation, $p$-adic, matrix, probability, special
function, and formal-verification approaches.  To avoid duplicating that work, the present
report takes a different route: it isolates the conjecture as a positivity-constrained
quadratic relation in the tensor square of the prime-logarithm space and studies the
\emph{rank} of that relation.

The main unconditional advance is a low-rank exclusion theorem.  Baker's homogeneous
linear-independence theorem implies that a nonzero rational quadratic form of rank at most
two cannot vanish at a tuple of $\Q$-linearly independent logarithms of positive algebraic
numbers.  Applied to Alaoglu--Erd\H{o}s, this proves that every hypothetical counterexample
must generate a genuinely rank-three or rank-four quadratic period relation.  We classify
the two possibilities exactly.  In multiplicative rank three, the relation is a smooth
rational conic with a distinguished rational point, and a rational column operation
symmetrizes the singular logarithm matrix.  Consequently a counterexample in this branch
would produce three multiplicatively independent positive algebraic numbers whose
logarithms form a geometric progression; more precisely, for
$\theta=\log 3/\log 2$ there would exist $\alpha>1$ algebraic such that
$\alpha$, $\alpha^\theta$, and $\alpha^{\theta^2}$ are all algebraic.  In multiplicative
rank four, the associated two-star quadratic has full rank four.

These reductions are rigorous, but they do not constitute an unconditional proof of the
conjecture.  The desired conclusion is an exact $2\times2$ case of the Four Exponentials
Conjecture, equivalently of the $2\times2$ Matrix Coefficient Conjecture.  The latter was
published in 2026 as a conjectural framework, with a key implication still left open.  The
report therefore states plainly which steps are theorems, which are equivalences, and which
single depth-two nonvanishing statement is still missing.  It also explains why recent
complete classifications of \emph{linear} relations among $1$-periods do not decide this
quadratic product relation.
\end{abstract}

\begin{auditbox}
No unconditional proof of the Alaoglu--Erd\H{o}s conjecture is presently obtained here, and
no such proof is claimed.  An assertion to the contrary would conceal an unproved instance
of Four Exponentials (or an equivalent matrix-coefficient statement).  What is proved is a
new, exact reduction: a counterexample can exist only through a rank-three or rank-four
quadratic relation among independent algebraic logarithms; all rank-one and rank-two
mechanisms are ruled out unconditionally by Baker's theorem.
\end{auditbox}

\tableofcontents
\newpage

\section{Scope, nonduplication, and theorem ledger}
\label{sec:scope}

\subsection{What this report does not repeat}

The attached unified report is already encyclopedic.  It includes, among many other themes,
the logarithmic determinant, the Four Exponentials reduction, exact conditional structure of
the solution monoid, prime-support and valuation lattices, rational-function rigidity,
finite-difference and interpolation arguments, determinant divisibility, $p$-adic and
tropical routes, matrix-coefficient formulations, point-counting strategies, probabilistic
models, special-function approaches, and a substantial Lean corpus.

Repeating those developments would add length without adding mathematical information.  The
present continuation therefore deliberately omits their proofs and develops four pieces that
are logically and technically different:
\begin{enumerate}[label=\textup{(\roman*)}]
\item a degree-two period map on the formal prime-exponent space;
\item a Baker-theoretic exclusion of every quadratic relation of rank at most two;
\item a hyperbolic pullback classification that identifies the only remaining ranks;
\item a rational symmetrization theorem for the multiplicatively dependent branch.
\end{enumerate}
The resulting report is self-contained for these claims and uses the broad earlier report
only as a reason to keep the scope focused.

\subsection{Ledger of claims}

For ease of audit, the logical status of the principal statements is recorded in
\Cref{tab:ledger}.

\begin{longtable}{@{}p{0.26\textwidth}p{0.25\textwidth}p{0.41\textwidth}@{}}
\toprule
Statement & Status & Content \\
\midrule
\endfirsthead
\toprule
Statement & Status & Content \\
\midrule
\endhead
Elementary reduction & Theorem & A nonintegral solution is positive, irrational, and in fact transcendental. \\
Four Exponentials reduction & Exact implication & Four Exponentials immediately rules out a nonintegral solution. \\
Restricted matrix coefficient & Exact equivalence at the relevant matrix & A counterexample gives a singular logarithm matrix with no vanishing rational matrix coefficient. \\
Prime-log two-star formulation & Exact equivalence & The conjecture is injectivity of a degree-two period map on a small positive cone. \\
Baker rank-two exclusion & Theorem & No nonzero rational quadratic form of rank $\le 2$ vanishes on independent algebraic logarithms. \\
Counterexample rank dichotomy & Theorem & The formal two-star quadratic of a counterexample has rank exactly $3$ or $4$. \\
Rank-three conic form & Theorem & The dependent branch is a nondegenerate split ternary conic with an explicit determinant and parametrization. \\
Symmetric logarithm normal form & Theorem & The dependent branch rationally symmetrizes to three independent logarithms in geometric progression. \\
Depth-two nonvanishing & Open & Exclude the remaining rank-three and rank-four two-star relations. \\
Alaoglu--Erd\H{o}s & Open in this report & Follows from Four Exponentials, Matrix Coefficient, Structural Rank, or Schanuel. \\
\bottomrule
\caption{Logical status of the main assertions.}
\label{tab:ledger}
\end{longtable}

\section{Elementary reductions and the singular logarithm matrix}
\label{sec:elementary}

Throughout, a \emph{candidate} is a triple $(x,M,N)$ with
\[
 x\in\R,\qquad M=2^x\in\Npos,\qquad N=3^x\in\Npos.
\]
A candidate is a \emph{counterexample} if $x\notin\Z$.

\subsection{The rational and algebraic cases}

\begin{proposition}[Elementary boundary]
\label{prop:elementary-boundary}
Let $(x,M,N)$ be a candidate.
\begin{enumerate}[label=\textup{(\alph*)}]
\item $x\ge 0$.
\item If $x\in\Q$, then $x\in\Nzero$.
\item If $x$ is algebraic, then $x\in\Nzero$.
\item Hence every counterexample has $x>0$ transcendental.
\end{enumerate}
\end{proposition}

\begin{proof}
If $x<0$, then $0<2^x<1$, contradicting $2^x\in\Z$.  Thus $x\ge0$.
Write a rational $x=a/b$ in lowest terms with $b>0$.  The identity
$2^a=M^b$ and unique factorization give $a=b\,\ord_2(M)$, so $b\mid a$.
Coprimality forces $b=1$, proving (b).  If $x$ is algebraic and irrational,
the Gelfond--Schneider theorem makes $2^x$ transcendental, contrary to
$M\in\Z$.  Thus algebraic $x$ is rational and (b) applies.  The last claim is immediate.
\end{proof}

The same unique-factorization argument shows that $\log 2$ and $\log 3$ are
$\Q$-linearly independent: a relation $a\log2+b\log3=0$ with integers $a,b$
would exponentiate to $2^a3^b=1$, forcing $a=b=0$.

\subsection{The rank-one real matrix and the rank-two formal problem}

Put
\[
 u=\log2,\qquad v=\log3,
\]
and form
\begin{equation}
 L=L(M,N)=
 \begin{pmatrix}
  \log2 & \log M\\
  \log3 & \log N
 \end{pmatrix}.
 \label{eq:log-matrix}
\end{equation}
For a candidate,
\begin{equation}
 L=
 \begin{pmatrix}u\\ v\end{pmatrix}
 \begin{pmatrix}1&x\end{pmatrix},
 \qquad
 \det L=u\log N-v\log M=0.
 \label{eq:outer-product}
\end{equation}
Thus the matrix has real rank one.  Its entries are logarithms of positive algebraic
numbers.  The problem is that their \emph{formal} dependence over $\Q$ need not have rank
one.

\begin{proposition}[Four Exponentials implication]
\label{prop:four-exp}
The Four Exponentials Conjecture implies the Alaoglu--Erd\H{o}s conjecture.
\end{proposition}

\begin{proof}
Suppose $x$ were a counterexample.  The pair $u,v$ is $\Q$-linearly independent, and so is
$1,x$.  Four Exponentials applied to the two rows $u,v$ and the two columns $1,x$ asserts
that at least one of
\[
 e^u=2,
 \quad e^{ux}=M,
 \quad e^v=3,
 \quad e^{vx}=N
\]
is transcendental.  All four are positive integers, a contradiction.
\end{proof}

\subsection{The exact matrix-coefficient obstruction}

For nonzero rational column vectors $w,z\in\Q^2$, 
\begin{equation}
 w^{T}Lz=(w_1u+w_2v)(z_1+xz_2).
 \label{eq:coefficient-factorization}
\end{equation}
This elementary factorization is the sharpest possible audit tool.

\begin{proposition}[Rational anisotropy]
\label{prop:anisotropy}
If $(x,M,N)$ is a counterexample, then
\[
 w^TLz\ne0
 \qquad\text{for every }w,z\in\Q^2\setminus\{0\}.
\]
\end{proposition}

\begin{proof}
The first factor in \eqref{eq:coefficient-factorization} is nonzero because $u,v$ are
$\Q$-linearly independent.  The second is nonzero because a relation
$z_1+xz_2=0$ with rational $z_1,z_2$, not both zero, would make $x$ rational.
\end{proof}

The $2\times2$ Matrix Coefficient Conjecture says that a singular matrix of logarithms of
algebraic numbers has such a vanishing rational matrix coefficient.  Therefore
\Cref{prop:anisotropy} says exactly that a counterexample would violate that conjecture.
Dasgupta and Kakde prove that the $2\times2$ Matrix Coefficient Conjecture is equivalent to
Four Exponentials and that the stronger Structural Rank Conjecture implies it; their paper
leaves open a key implication in the proposed proof strategy \cite{DK2026}.  This is not a
historical side note: it identifies the precise contemporary frontier occupied by
\eqref{eq:log-matrix}.

\begin{frontierbox}
A proof that starts only from the facts ``$L$ is singular'' and ``the entries of $L$ are
logarithms of algebraic numbers'' must supply, in this special case, the missing conclusion
of the Matrix Coefficient Conjecture.  The factorization
\eqref{eq:coefficient-factorization} then finishes the argument instantly.
\end{frontierbox}

\section{The prime-exponent space and a depth-two period map}
\label{sec:depth-two}

\subsection{Formal prime logarithms}

Let $\mathcal P$ be the set of primes and let
\[
 V=\bigoplus_{p\in\mathcal P}\Q e_p
\]
be the rational vector space with one basis vector for each prime.  For a positive rational
number $r$, define its exponent vector
\[
 \nu(r)=\sum_{p\in\mathcal P}\ord_p(r)e_p\in V.
\]
Define a linear evaluation map
\begin{equation}
 \ellmap:V\longrightarrow\R,
 \qquad
 \ellmap(e_p)=\log p.
 \label{eq:ell-map}
\end{equation}
Then $\ellmap(\nu(r))=\log r$.

\begin{lemma}[Unique-factorization injectivity]
\label{lem:ell-injective}
The map $\ellmap$ is injective.
\end{lemma}

\begin{proof}
If $\sum_p c_p\log p=0$ with finitely many $c_p\in\Q$, clear denominators to obtain
integers $n_p$ with $\sum_p n_p\log p=0$.  Exponentiation gives
$\prod_p p^{n_p}=1$, and unique factorization forces every $n_p=0$.
\end{proof}

The first symmetric power $V$ records linear logarithmic relations.  The determinant
condition lies in the second symmetric power.  Extend \eqref{eq:ell-map} multiplicatively to
\begin{equation}
 \pi_2:\Sym^2(V)\longrightarrow\R,
 \qquad
 \pi_2(rs)=\ellmap(r)\ellmap(s)
 \quad(r,s\in V).
 \label{eq:pi2}
\end{equation}
This is the degree-two, or ``depth-two,'' period map used in this report.  The word
``depth'' is local terminology: it simply distinguishes products of logarithms from linear
relations among logarithms.

\subsection{The two-star tensor}

Let
\[
 a=\nu(M),\qquad b=\nu(N).
\]
Define
\begin{equation}
 q_{a,b}=e_2b-e_3a\in\Sym^2(V).
 \label{eq:two-star}
\end{equation}
It has a graph-theoretic shape consisting of two stars, one centered at $2$ and one centered
at $3$, with opposite signs.  Evaluation gives
\begin{equation}
 \pi_2(q_{a,b})
 =\log2\log N-\log3\log M.
 \label{eq:two-star-eval}
\end{equation}
Thus every candidate satisfies $\pi_2(q_{a,b})=0$.

To describe the formal zero, write
\[
 a=\sum_p a_pe_p,
 \qquad
 b=\sum_p b_pe_p.
\]
The coefficient of $e_2e_p$ for $p\notin\{2,3\}$ is $b_p$, the coefficient of
$e_3e_p$ is $-a_p$, the coefficients of $e_2^2,e_3^2$ are $b_2,-a_3$, and the coefficient
of $e_2e_3$ is $b_3-a_2$.

\begin{lemma}[Formal zero classification]
\label{lem:formal-zero}
For $a,b\in\bigoplus_p\Nzero e_p$, one has $q_{a,b}=0$ in $\Sym^2(V)$ if and only if
there is $m\in\Nzero$ such that
\[
 a=me_2,
 \qquad
 b=me_3.
\]
\end{lemma}

\begin{proof}
If $q_{a,b}=0$, the coefficients just listed show that $b_p=0$ for $p\ne3$, that
$a_p=0$ for $p\ne2$, and that $b_3=a_2$.  Calling the common value $m$ gives the stated
form.  The converse is immediate.
\end{proof}

Define the positive two-star cone
\begin{equation}
 \mathcal C_{2,3}
 =\{e_2b-e_3a:a,b\in\textstyle\bigoplus_p\Nzero e_p\}
 \subset\Sym^2(V).
 \label{eq:two-star-cone}
\end{equation}

\begin{theorem}[Exact depth-two formulation]
\label{thm:depth-two-equivalence}
The following are equivalent.
\begin{enumerate}[label=\textup{(\roman*)}]
\item For every real $x$, $2^x,3^x\in\Z$ implies $x\in\Z$.
\item $\ker\pi_2\cap\mathcal C_{2,3}=\{0\}$.
\item For all nonnegative integral exponent vectors $a,b$, the equality
\[
 \log2\,\ellmap(b)=\log3\,\ellmap(a)
\]
forces $a=me_2$, $b=me_3$ for some $m\in\Nzero$.
\end{enumerate}
\end{theorem}

\begin{proof}
The equivalence of (ii) and (iii) is \Cref{lem:formal-zero}.  Suppose (ii), and let
$(x,M,N)$ be a candidate.  By \eqref{eq:two-star-eval}, $q_{a,b}\in\ker\pi_2$; hence it is
formally zero, so $M=2^m$, $N=3^m$.  Since $2^x=2^m$, one has $x=m$.
Conversely, if (ii) fails, choose nonnegative integral $a,b$ with nonzero $q_{a,b}$ and
$\pi_2(q_{a,b})=0$.  Put $M=\prod p^{a_p}$ and $N=\prod p^{b_p}$.  Then
\[
 x=\frac{\log M}{\log2}=\frac{\log N}{\log3}
\]
gives $2^x=M$ and $3^x=N$, while \Cref{lem:formal-zero} shows $x\notin\Z$.
\end{proof}

\begin{frontierbox}
The conjecture is not asking for all quadratic relations among prime logarithms to vanish.
It asks only for injectivity of $\pi_2$ on the extremely sparse, positivity-constrained cone
$\mathcal C_{2,3}$.  This is much weaker than algebraic independence of prime logarithms,
but it is genuinely quadratic and therefore lies beyond Baker's linear theorem.
\end{frontierbox}

\section{Baker's theorem eliminates quadratic rank at most two}
\label{sec:baker-rank}

\subsection{Rank of a quadratic tensor}

If $U$ is a finite-dimensional rational vector space and $q\in\Sym^2(U)$, contraction makes
$q$ a symmetric linear map
\[
 q^\sharp:U^*\longrightarrow U.
\]
Its rank is the \emph{rank of $q$}.  In coordinates this is the rank of any symmetric
coefficient matrix for the quadratic form, and is independent of the harmless convention
of putting factors $2$ on mixed terms.

We use the following qualitative form of Baker's theorem.

\begin{theorem}[Baker, homogeneous form]
\label{thm:baker}
Let $\lambda_1,\ldots,\lambda_n$ be logarithms of nonzero algebraic numbers.  If they are
linearly independent over $\Q$, then they are linearly independent over $\Qbar$.
\end{theorem}

For positive algebraic numbers, the real logarithms remove all branch issues.

\begin{lemma}[Low-rank factorization]
\label{lem:low-rank-factor}
Let $K$ be an algebraically closed field of characteristic different from $2$.  Every
nonzero quadratic form over $K$ of rank one or two is a product of two nonzero linear forms
over $K$.
\end{lemma}

\begin{proof}
After an invertible linear change of variables, a rank-one form is $X_1^2$.  A rank-two form
is $X_1^2+X_2^2$, which factors over the algebraically closed field as
$(X_1+iX_2)(X_1-iX_2)$ after choosing $i^2=-1$.
\end{proof}

\begin{theorem}[Baker rank-two exclusion]
\label{thm:baker-rank-two}
Let $\lambda_1,\ldots,\lambda_n$ be $\Q$-linearly independent logarithms of positive
algebraic numbers.  If $0\ne q\in\Q[X_1,\ldots,X_n]$ is homogeneous quadratic and
$\rank(q)\le2$, then
\[
 q(\lambda_1,\ldots,\lambda_n)\ne0.
\]
\end{theorem}

\begin{proof}
By \Cref{lem:low-rank-factor}, over $\Qbar$ one may write $q=L_1L_2$, where $L_1,L_2$ are
nonzero algebraic linear forms; in rank one they may coincide.  If $q(\lambda)=0$, then
$L_1(\lambda)=0$ or $L_2(\lambda)=0$.  This is a nontrivial $\Qbar$-linear relation among
the $\lambda_i$, contradicting \Cref{thm:baker}.
\end{proof}

\begin{resultbox}
Theorem~\ref{thm:baker-rank-two} is a genuine quadratic consequence of Baker's linear
independence theorem.  It does not prove that arbitrary quadratic forms are nonzero; it
identifies an exact algebraic boundary.  Rank one and rank two factor into linear forms and
are therefore impossible.  Rank three is the first case in which factorization no longer
reduces the problem to Baker.
\end{resultbox}

\subsection{Application to the two-star tensor}

Let $U=\Span_{\Q}\{e_2,e_3,a,b\}$.  Choose a rational basis
$f_1,\ldots,f_r$ of $U$.  The numbers $\ellmap(f_i)$ are logarithms of positive algebraic
numbers: after clearing denominators, each $f_i$ is an integral combination of prime basis
vectors, and division by the common denominator corresponds to a positive algebraic root.
They are $\Q$-linearly independent by \Cref{lem:ell-injective}.

\begin{corollary}[First unconditional rank restriction]
\label{cor:first-rank-restriction}
If $(x,M,N)$ is a counterexample, then the quadratic tensor $q_{a,b}$ has rank at least
three.
\end{corollary}

\begin{proof}
It is nonzero by \Cref{lem:formal-zero}, and its evaluation at the independent algebraic
logarithms $\ellmap(f_i)$ vanishes by \eqref{eq:two-star-eval}.  Theorem
\ref{thm:baker-rank-two} excludes rank at most two.
\end{proof}

This restriction is stronger than the familiar statement that a counterexample must involve
some prime outside $\{2,3\}$.  It says that the entire quadratic tensor cannot secretly
collapse to a product of two algebraic linear forms.

\section{Hyperbolic pullback and the exact rank classification}
\label{sec:hyperbolic}

\subsection{A four-vector theorem}

The special shape $eb-fa$ admits a complete linear-algebra analysis.

\begin{theorem}[Hyperbolic pullback classification]
\label{thm:hyperbolic-classification}
Let $V$ be a rational vector space, let $e,f\in V$ be linearly independent, and let
$a,b\in V$.  Put
\[
 q=eb-fa\in\Sym^2(V),
 \qquad
 r=\dim_{\Q}\Span\{e,b,f,a\}.
\]
Then:
\begin{enumerate}[label=\textup{(\roman*)}]
\item if $r=4$, then $\rank(q)=4$;
\item if $r\le2$, then $\rank(q)\le2$;
\item if $r=3$, let
\begin{equation}
 c_1e+c_2b+c_3f+c_4a=0
 \label{eq:unique-relation}
\end{equation}
be a nonzero relation, unique up to scale.  Then
\[
 \rank(q)=
 \begin{cases}
 2,&c_1c_2-c_3c_4=0,\\
 3,&c_1c_2-c_3c_4\ne0.
 \end{cases}
\]
\end{enumerate}
\end{theorem}

\begin{proof}
Let $E=\Q^4$ with basis $z_1,z_2,z_3,z_4$ and let
$\phi:E\to\Span\{e,b,f,a\}$ send the basis to $e,b,f,a$.  The tensor $q$ is the image of
\[
 q_0=z_1z_2-z_3z_4.
\]
As a bilinear form on $E^*$, twice $q_0$ has matrix
\[
 G=
 \begin{pmatrix}
 0&1&0&0\\
 1&0&0&0\\
 0&0&0&-1\\
 0&0&-1&0
 \end{pmatrix},
 \qquad G^{-1}=G.
\]
The bilinear form defined by $q$ is the restriction of this hyperbolic form to the
$r$-dimensional subspace $\operatorname{im}(\phi^*)\subset E^*$.

If $r=4$, the map $\phi$ is an isomorphism and the rank is $4$.  If $r\le2$, restriction to
an at most two-dimensional space has rank at most $2$.

Now suppose $r=3$.  The kernel of $\phi$ is the line spanned by
$c=(c_1,c_2,c_3,c_4)^T$, and
\[
 \operatorname{im}(\phi^*)=c^\perp
 =\{z\in E^*:c^Tz=0\}.
\]
The restriction of $G$ to $c^\perp$ is degenerate precisely when its $G$-normal
$G^{-1}c=Gc$ also lies in $c^\perp$, equivalently when
\[
 c^TG^{-1}c=2(c_1c_2-c_3c_4)=0.
\]
In that case the radical is one-dimensional and the restricted rank is $2$; otherwise the
restriction is nondegenerate of rank $3$.
\end{proof}

The scalar
\begin{equation}
 \Delta_{\mathrm{tan}}(c)=c_1c_2-c_3c_4
 \label{eq:tangent-invariant}
\end{equation}
can be viewed as a tangency invariant: the relation hyperplane is tangent to the hyperbolic
quadric $z_1z_2-z_3z_4=0$ exactly when $\Delta_{\mathrm{tan}}(c)=0$.

\subsection{Multiplicative rank of the four integers}

Let
\begin{equation}
 r_{\times}(2,3,M,N)
 =\dim_{\Q}\Span\{e_2,e_3,a,b\}.
 \label{eq:mult-rank}
\end{equation}
This is the rational multiplicative rank of the subgroup generated by the four positive
integers.  It is either $2$, $3$, or $4$ because $e_2,e_3$ are independent.

\begin{theorem}[Counterexample rank theorem]
\label{thm:counterexample-rank}
For a counterexample $(x,M,N)$,
\[
 r_{\times}(2,3,M,N)\in\{3,4\},
 \qquad
 \rank(q_{a,b})=r_{\times}(2,3,M,N).
\]
More explicitly:
\begin{enumerate}[label=\textup{(\alph*)}]
\item in multiplicative rank four, $q_{a,b}$ has rank four;
\item in multiplicative rank three, $q_{a,b}$ has rank three;
\item multiplicative rank two is impossible.
\end{enumerate}
\end{theorem}

\begin{proof}
If the multiplicative rank is four, \Cref{thm:hyperbolic-classification} gives rank four.
If it is at most two, the same theorem gives quadratic rank at most two, contradicting
\Cref{cor:first-rank-restriction}.  In multiplicative rank three the quadratic rank is two or
three by \Cref{thm:hyperbolic-classification}, and Baker excludes rank two.
\end{proof}

\begin{resultbox}
A hypothetical counterexample cannot be explained by any hidden low-dimensional
factorization.  Its formal quadratic relation has exactly the same rank as the
multiplicative group generated by $2,3,M,N$, and that common rank is $3$ or $4$.
\end{resultbox}

\subsection{The unique relation in rank three}

Suppose the multiplicative rank is three.  Then there is a unique rational relation, up to
scale,
\begin{equation}
 M^{\alpha}N^{\beta}2^{\gamma}3^{\delta}=1,
 \qquad
 (\alpha,\beta,\gamma,\delta)\in\Z^4\setminus\{0\}.
 \label{eq:multiplicative-relation}
\end{equation}
In exponent space,
\[
 \gamma e_2+\beta b+\delta e_3+\alpha a=0.
\]
Thus the vector $c$ in \eqref{eq:unique-relation}, for the ordering
$(e,b,f,a)=(e_2,b,e_3,a)$, is
\[
 c=(\gamma,\beta,\delta,\alpha).
\]
The tangent invariant is
\begin{equation}
 \Delta_{\mathrm{tan}}=\gamma\beta-\delta\alpha.
 \label{eq:det-invariant}
\end{equation}

\begin{corollary}[Non-tangency]
\label{cor:non-tangency}
For a rank-three counterexample,
\[
 \gamma\beta-\delta\alpha\ne0.
\]
\end{corollary}

\begin{proof}
If it vanished, \Cref{thm:hyperbolic-classification} would give
$\rank(q_{a,b})=2$, contradicting \Cref{thm:counterexample-rank}.
\end{proof}

There is also a direct interpretation.  Taking logarithms in
\eqref{eq:multiplicative-relation} and substituting
$\log M=x\log2$, $\log N=x\log3$ gives
\[
 x(\alpha+\beta\theta)+(\gamma+\delta\theta)=0,
 \qquad \theta=\frac{\log3}{\log2}.
\]
If $\gamma\beta-\delta\alpha=0$, numerator and denominator are rationally proportional,
forcing $x\in\Q$, which has already been excluded.  The rank calculation shows that this
familiar fractional-linear determinant is exactly the dual hyperbolic tangency invariant.

\section{The rank-three conic and its explicit arithmetic normal form}
\label{sec:rank-three-conic}

\subsection{A common external prime direction}

Decompose
\[
 V=\Span\{e_2,e_3\}\oplus V_{\mathrm{ext}},
 \qquad
 V_{\mathrm{ext}}=\bigoplus_{p\ne2,3}\Q e_p.
\]
Write $a_{\mathrm{ext}},b_{\mathrm{ext}}$ for the projections of $a,b$.  Then
\begin{equation}
 r_{\times}(2,3,M,N)
 =2+\dim_{\Q}\Span\{a_{\mathrm{ext}},b_{\mathrm{ext}}\}.
 \label{eq:external-rank}
\end{equation}
Consequently the rank-three branch is exactly the case in which the two external exponent
vectors span one rational line.

After choosing a primitive nonnegative integral vector $h\in V_{\mathrm{ext}}$, there are
$u,v\in\Nzero$, not both zero, and nonnegative integers
$A,B,C,D$ such that
\begin{equation}
 \begin{aligned}
 M&=2^A3^B R^u,\\
 N&=2^C3^D R^v,
 \end{aligned}
 \qquad
 R=\prod_{p\ne2,3}p^{h_p}>1.
 \label{eq:common-external-core}
\end{equation}
One of $u,v$ may vanish, corresponding to one output having no external prime factor.

Put
\[
 X=\log2,
 \qquad Y=\log3,
 \qquad Z=\log R.
\]
The determinant relation becomes
\begin{equation}
 Q(X,Y,Z)
 =CX^2+(D-A)XY-BY^2+vXZ-uYZ=0.
 \label{eq:ternary-q}
\end{equation}
The three logarithms $X,Y,Z$ are $\Q$-linearly independent by unique factorization.

\subsection{The determinant certificate}

Twice the symmetric matrix of \eqref{eq:ternary-q} is
\begin{equation}
 H=
 \begin{pmatrix}
 2C&D-A&v\\
 D-A&-2B&-u\\
 v&-u&0
 \end{pmatrix}.
 \label{eq:H-matrix}
\end{equation}
A direct calculation gives
\begin{equation}
 \det H
 =2\bigl((A-D)uv+Bv^2-Cu^2\bigr).
 \label{eq:H-determinant}
\end{equation}

\begin{proposition}[Rank-three boundary certificate]
\label{prop:boundary-certificate}
A rank-three counterexample represented as in \eqref{eq:common-external-core} must satisfy
\begin{equation}
 (A-D)uv+Bv^2-Cu^2\ne0.
 \label{eq:boundary-nonzero}
\end{equation}
Equivalently, the ternary quadratic form \eqref{eq:ternary-q} is nondegenerate.
\end{proposition}

\begin{proof}
The rank of $Q$ is the rank of $H$.  Theorem~\ref{thm:counterexample-rank} says that this rank
is three, so the determinant cannot vanish.
\end{proof}

The form $Q$ nevertheless has the rational isotropic point $[0:0:1]$.  Its conic is therefore
split and rationally parametrizable.  The point is smooth because the gradient there is
$(v,-u,0)$, which is nonzero.

\subsection{Parametrization at the prime-log point}

Set
\[
 \theta=\frac{Y}{X}=\log_2 3,
 \qquad
 \rho=\frac{Z}{X}=\log_2 R.
\]
Dividing \eqref{eq:ternary-q} by $X^2$ gives
\begin{equation}
 C+(D-A)\theta-B\theta^2+(v-u\theta)\rho=0.
 \label{eq:conic-affine}
\end{equation}
Since $\theta$ is irrational and $u,v$ are rational, not both zero,
$v-u\theta\ne0$.  Hence
\begin{equation}
 \rho=
 \frac{B\theta^2-(D-A)\theta-C}{v-u\theta}.
 \label{eq:rho-rational-function}
\end{equation}

\begin{proposition}[Explicit rank-three obstruction]
\label{prop:rank-three-obstruction}
Every rank-three counterexample produces integers $A,B,C,D,u,v$ satisfying
\eqref{eq:boundary-nonzero} and an integer $R>1$, supported away from $2$ and $3$, for which
\eqref{eq:rho-rational-function} holds with
$\theta=\log_2 3$ and $\rho=\log_2 R$.
Conversely, any such data with \eqref{eq:conic-affine} gives a vanishing two-star quadratic
of multiplicative rank three; it gives an Alaoglu--Erd\H{o}s counterexample exactly when the
corresponding ratios
\[
 \frac{\log(2^A3^BR^u)}{\log2}
 \quad\text{and}\quad
 \frac{\log(2^C3^DR^v)}{\log3}
\]
are equal and nonintegral, which is precisely \eqref{eq:conic-affine} plus nonintegrality.
\end{proposition}

The formula is a useful compression of the dependent branch.  It does not prove
nonexistence: proving that $2^{F(\log_2 3)}$ is transcendental for the relevant nonlinear
rational functions $F$ is itself beyond current theorems and contains a Four Exponentials
type difficulty.

\section{Rational symmetrization and a geometric progression of logarithms}
\label{sec:symmetrization}

The conic description is coordinate-dependent.  A more invariant normal form comes directly
from the $2\times2$ logarithm matrix.

\subsection{Symmetrizing a rank-three singular logarithm matrix}

For $C,L\in M_2(\R)$ use the Frobenius pairing
\[
 \langle C,L\rangle_F=\operatorname{tr}(C^TL)=\sum_{i,j}C_{ij}L_{ij}.
\]
Let
\[
 J=\begin{pmatrix}0&1\\-1&0\end{pmatrix}.
\]

\begin{theorem}[Rational symmetrization]
\label{thm:rational-symmetrization}
Let $L\in M_2(\R)$ be nonzero and singular.  Assume that its four entries span a
three-dimensional $\Q$-vector space and that the unique rational relation among them has an
invertible coefficient matrix $C\in M_2(\Q)$:
\[
 \langle C,L\rangle_F=0,
 \qquad \det C\ne0.
\]
Then there is $Q\in\mathrm{GL}_2(\Q)$ such that $S=LQ$ is symmetric.  If, in addition,
$L=r(1,x)$ is an outer product with $r=(r_1,r_2)^T$, then
\begin{equation}
 S=\kappa
 \begin{pmatrix}
 r_1^2&r_1r_2\\
 r_1r_2&r_2^2
 \end{pmatrix}
 \label{eq:symmetric-rank-one-form}
\end{equation}
for some nonzero real $\kappa$.
\end{theorem}

\begin{proof}
Under the right change $S=LQ$, the relation matrix becomes
$C'=CQ^{-T}$ because
\[
 \langle C,L\rangle_F
 =\langle CQ^{-T},LQ\rangle_F.
\]
Choose $Q$ by $Q^{-T}=C^{-1}J$.  Then $C'=J$, and
\[
 0=\langle J,S\rangle_F=S_{12}-S_{21},
\]
so $S$ is symmetric.

For the second assertion, write $(1,x)Q=d^T$.  Then $S=rd^T$.  Symmetry says
$r_1d_2=r_2d_1$.  Since $r\ne0$ and $d\ne0$, the vectors $d$ and $r$ are proportional:
$d=\kappa r$.  This gives \eqref{eq:symmetric-rank-one-form}.
\end{proof}

\subsection{Application to a dependent counterexample}

For the logarithm matrix \eqref{eq:log-matrix}, a multiplicative relation
\eqref{eq:multiplicative-relation} is
\[
 \left\langle
 \begin{pmatrix}\gamma&\alpha\\\delta&\beta\end{pmatrix},
 L\right\rangle_F=0.
\]
Its determinant is $\gamma\beta-\delta\alpha$, which is nonzero by
\Cref{cor:non-tangency}.  Therefore \Cref{thm:rational-symmetrization} applies with
$r=(u,v)^T=(\log2,\log3)^T$.

\begin{theorem}[Geometric-progression normal form]
\label{thm:gp-normal-form}
Suppose $(x,M,N)$ is a multiplicative-rank-three counterexample.  Then there exists
$\kappa>0$ such that the three numbers
\begin{equation}
 \lambda_0=\kappa(\log2)^2,
 \qquad
 \lambda_1=\kappa\log2\log3,
 \qquad
 \lambda_2=\kappa(\log3)^2
 \label{eq:lambda-progression}
\end{equation}
are logarithms of positive algebraic numbers, are $\Q$-linearly independent, and satisfy
\begin{equation}
 \lambda_1^2=\lambda_0\lambda_2,
 \qquad
 \frac{\lambda_1}{\lambda_0}
 =\frac{\lambda_2}{\lambda_1}
 =\theta=\frac{\log3}{\log2}.
 \label{eq:gp-logs}
\end{equation}
Equivalently, there exists a positive algebraic number $\alpha>1$ such that
\begin{equation}
 \alpha,\qquad \alpha^\theta,\qquad \alpha^{\theta^2}
 \quad\text{are all algebraic and multiplicatively independent.}
 \label{eq:theta-orbit}
\end{equation}
\end{theorem}

\begin{proof}
The preceding discussion gives $Q\in\mathrm{GL}_2(\Q)$ and nonzero $\kappa$ with
\[
 LQ=\kappa
 \begin{pmatrix}u^2&uv\\uv&v^2\end{pmatrix}.
\]
Replacing $Q$ by $-Q$ if necessary makes $\kappa>0$.  Each entry of $LQ$ is a rational
linear combination of logarithms of positive integers.  Its exponential is therefore a
positive algebraic number: rational coefficients correspond to products, quotients, and
positive real algebraic roots.  This proves algebraicity of the exponentials of the three
numbers in \eqref{eq:lambda-progression}.

Right multiplication by an invertible rational matrix preserves the $\Q$-span of all matrix
entries.  In multiplicative rank three, the entries of $L$ span dimension three.  The
symmetric matrix $LQ$ has only the three displayed distinct entries, so they are
$\Q$-linearly independent.  The geometric-progression identities are immediate.

Set $\alpha=e^{\lambda_0}>1$.  Then
$e^{\lambda_1}=\alpha^\theta$ and
$e^{\lambda_2}=\alpha^{\theta^2}$ are algebraic.  Independence of the logarithms is exactly
multiplicative independence of these three algebraic numbers.
\end{proof}

\begin{corollary}[The common ratio is necessarily transcendental]
\label{cor:theta-transcendental-again}
In the situation of \Cref{thm:gp-normal-form}, the common ratio $\theta$ cannot be algebraic.
\end{corollary}

\begin{proof}
It is not rational because $\log2$ and $\log3$ are $\Q$-linearly independent.  If it were
algebraic irrational, Gelfond--Schneider applied to the algebraic base $\alpha\ne0,1$ would
make $\alpha^\theta$ transcendental, contradicting \eqref{eq:theta-orbit}.
\end{proof}

Of course the transcendence of $\theta=\log_2 3$ is already classical by the same theorem,
because $2^\theta=3$.  The point of \Cref{thm:gp-normal-form} is not to reprove that fact; it
is to show that the entire multiplicatively dependent branch of Alaoglu--Erd\H{o}s has a
single symmetric shape.

\begin{resultbox}
The rank-three branch is forced into a three-term algebraic orbit
$\alpha,\alpha^\theta,\alpha^{\theta^2}$.  This is a smaller and more symmetric obstruction
than the original four entries $2,3,M,N$, but it is still a Four Exponentials type
obstruction: current transcendence theorems do not exclude a transcendental exponent
$\theta$ from having two consecutive algebraic power iterates.
\end{resultbox}

\section{Why the complete theory of linear periods does not settle the problem}
\label{sec:period-audit}

\subsection{Degree one versus degree two}

Each $\log\alpha$ for algebraic $\alpha$ is a classical $1$-period.  Baker's theorem is a
complete rigidity statement for a fundamental class of \emph{linear} relations.  The modern
work of Huber and W\"ustholz gives a conceptual classification of linear relations among
$1$-periods via $1$-motives \cite{HW2022}; recent work makes this picture more complete and,
in important formulations, effective \cite{HuberKalck2026,SertozOuaknineWorrell2025}.

The Alaoglu--Erd\H{o}s determinant is not a linear relation among the four logarithms.  It is
\[
 (\log2)(\log N)-(\log3)(\log M)=0,
\]
a linear relation between \emph{products} of logarithms.  In the notation of
\Cref{sec:depth-two}, it lies in $\Sym^2(V)$ and asks about $\ker\pi_2$, not
$\ker\ellmap$.  Products of periods are again periods, but a complete description of their
quadratic relations would be a much deeper period statement.

\begin{proposition}[Linear-period completeness cannot be applied verbatim]
\label{prop:linear-period-limit}
Any theorem whose conclusion controls only $\Qbar$-linear relations among the individual
numbers
\[
 \log2,\ \log3,\ \log M,\ \log N
\]
does not by itself exclude the determinant relation
\[
 \log2\log N-\log3\log M=0.
\]
\end{proposition}

\begin{proof}
A counterexample may have the four logarithms $\Qbar$-linearly independent in multiplicative
rank four, by Baker's theorem, while still satisfying the displayed quadratic relation.
Linear independence in a vector space does not imply linear independence of all products in
its symmetric square.  In multiplicative rank three there is one rational linear relation,
but after choosing a basis of three independent logs the determinant becomes a nondegenerate
rank-three quadratic form, again outside the scope of a theorem about linear relations.
\end{proof}

\subsection{A torus-homomorphism no-go lemma}

The analytic subgroup theorem is extraordinarily powerful for linear logarithmic relations.
It is important, however, not to assign to it a nonlinear conclusion that its functorial
input does not encode.

\begin{lemma}[Differentials of torus homomorphisms are integral-linear]
\label{lem:torus-linear}
Every algebraic-group homomorphism
$\varphi:\Gm^r\to\Gm^s$ has the form
\[
 \varphi(z_1,\ldots,z_r)_j=\prod_{i=1}^r z_i^{n_{ji}}
 \qquad(n_{ji}\in\Z),
\]
and its differential at the identity sends logarithmic coordinates
$\lambda\in\C^r$ to $N\lambda$, where $N=(n_{ji})$.
\end{lemma}

\begin{proof}
The character group of $\Gm^r$ is $\Z^r$.  Pullback of each coordinate character of
$\Gm^s$ is therefore an integral monomial character on $\Gm^r$, giving the displayed form.
Differentiating the monomials at the identity gives the integer matrix $N$.
\end{proof}

\begin{corollary}[No automatic linearization of the determinant]
\label{cor:no-automatic-linearization}
Constructions using only products of algebraic tori, algebraic subgroups, quotients, and
algebraic homomorphisms transform logarithmic coordinates by rational or integral linear
maps.  They do not turn the quadratic tensor
$e_2b-e_3a$ into a linear logarithmic relation unless an additional genuinely degree-two
identity is supplied.
\end{corollary}

This is a scope statement, not a claim that motivic methods can never solve the problem.
Tensor products of motives and stronger period conjectures are precisely mechanisms by which
degree-two information might enter.  The point is that the already-complete degree-one
classification is not the missing theorem.

\subsection{Conditional period consequences}

The following standard conjectural principles all imply the desired nonvanishing:
\begin{itemize}
\item Schanuel's conjecture implies sufficient algebraic independence of independent
logarithms;
\item the Structural Rank Conjecture implies the Matrix Coefficient Conjecture;
\item the Matrix Coefficient Conjecture in size two is equivalent to Four Exponentials;
\item injectivity of $\pi_2$ on all of $\Sym^2(V)$ is far stronger than needed but immediately
settles the positive two-star cone.
\end{itemize}
The report does not assume any of these conjectures in an unconditional theorem.

\section{The 2026 matrix frontier in exact form}
\label{sec:matrix-frontier}

\subsection{Three matrix statements}

Let $\mathscr L$ be the $\Q$-vector space of complex logarithms of algebraic numbers.
For a matrix $A\in M_{m,n}(\mathscr L)$, write
\[
 A=\sum_{k=1}^r A_k\lambda_k,
\]
where the $A_k$ have rational entries and the $\lambda_k$ are $\Q$-linearly independent.
The relevant conjectures are:

\begin{description}[leftmargin=3.4cm,style=nextline]
\item[Structural Rank.] The numerical rank of $A$ equals the rank of the generic matrix
$\sum A_kX_k$ over $\Q(X_1,\ldots,X_r)$.
\item[Matrix Coefficient.] If a square $A$ is singular, there are nonzero rational vectors
$w,z$ with $w^TAz=0$.
\item[Four Exponentials.] If two row parameters and two column parameters are each
$\Q$-linearly independent, at least one of the four corresponding exponentials is
transcendental.
\end{description}

Dasgupta and Kakde establish
\[
 \SRC\Longrightarrow\MCC,
 \qquad
 \MCC\Longleftrightarrow\FE
 \quad\text{in size }2,
\]
and develop an auxiliary-polynomial strategy, while explicitly leaving a needed implication
open \cite{DK2026}.  The paper appeared in \emph{Advances in Mathematics} in March 2026.

\subsection{Specialization to Alaoglu--Erd\H{o}s}

For a nonintegral candidate, $L$ in \eqref{eq:log-matrix} is singular.  If $\MCC$ holds,
there are nonzero $w,z\in\Q^2$ with $w^TLz=0$.  Proposition
\ref{prop:anisotropy} says this is impossible.  Thus:

\begin{theorem}[Conditional one-line proof]
\label{thm:conditional-one-line}
Each of the following implies the Alaoglu--Erd\H{o}s conjecture:
\[
 \text{Schanuel}
 \Longrightarrow \SRC
 \Longrightarrow \MCC
 \Longleftrightarrow \FE
 \Longrightarrow \AEconj.
\]
For the final implication, the complete proof is the factorization
\[
 w^TLz=(w_1\log2+w_2\log3)(z_1+xz_2).
\]
\end{theorem}

The difficult part is not the final line; it is obtaining the vanishing matrix coefficient
from singularity.

\subsection{Why Six Exponentials stops one row or column short}

The Six Exponentials Theorem handles a $3\times2$ or $2\times3$ grid.  A counterexample
supplies only the $2\times2$ algebraic grid
\[
 \begin{array}{c|cc}
  &1&x\\ \hline
 \log2&2&M\\
 \log3&3&N
 \end{array}.
\]
To add a third row $\log R$, one would need algebraicity of $R^x$; to add a third column
$y$, one would need algebraicity of both $2^y$ and $3^y$.  Neither is available from the
hypotheses.  In the rank-three normal form, Six Exponentials similarly sees
$\alpha,\alpha^\theta,\alpha^{\theta^2}$ but requires one additional algebraic entry to force
a contradiction.

\section{Audit of tempting but invalid proof routes}
\label{sec:false-proofs}

This section records failure modes because the same hidden assumptions repeatedly reappear in
attempts at the problem.

\subsection{Taking prime valuations of a real power}

From $M=2^x$ one cannot write
$\ord_p(M)=x\ord_p(2)$ unless $x$ is already known to be an integer in an algebraic setting
where such a valuation identity is defined.  The real exponential $2^x$ is defined
analytically; valuation is applied only after the output happens to be an integer.  Extending
$\ord_p(a^x)=x\ord_p(a)$ to arbitrary real $x$ assumes precisely the arithmetic structure at
issue.

\subsection{Applying Baker to a product relation}

Baker excludes nontrivial algebraic \emph{linear} combinations
$\sum\beta_i\log\alpha_i$.  The determinant
\[
 \log2\log N-\log3\log M
\]
has logarithms as coefficients of other logarithms, and those coefficients are generally
transcendental.  It is not a Baker linear form.  The valid consequence of Baker is exactly
\Cref{thm:baker-rank-two}: factorable low-rank quadratics are excluded, while irreducible
rank-three and rank-four quadratics remain.

\subsection{Treating the exponent as algebraic}

Gelfond--Schneider proves the conjecture when $x$ is algebraic.  A counterexample, however,
would necessarily have transcendental $x$.  Any proof that invokes Gelfond--Schneider on
$2^x$ without first proving algebraicity of $x$ is circular.

\subsection{Manufacturing a third base}

The Six Exponentials Theorem would solve the problem if there were a third
multiplicatively independent algebraic base $R$ with $R^x$ algebraic.  The hypotheses give
this only for $R$ in the group generated by $2$ and $3$, which adds no independent row.
Choosing an external prime and asserting algebraicity of its $x$-th power is an extra
assumption.

\subsection{A noncanonical \texorpdfstring{$p$}{p}-adic exponent}

The real identity $2^x=M$ does not canonically define a $p$-adic number $x$ or a compatible
$p$-adic power $2^x$.  One may form $p$-adic logarithms of suitable algebraic numbers after
choosing domains and branches, but transferring the same real transcendental exponent into
$\Q_p$ requires new data.  A proof that silently identifies the real and $p$-adic exponents
has changed the problem.

\subsection{Integer gaps and rational approximation}

Let $p/q$ be a convergent to a hypothetical irrational $x$.  Then
\[
 |qx-p|<\frac1q.
\]
The integer gap $|M^q-2^p|\ge1$ gives only an exponentially small lower bound on
$|qx-p|$.  A Baker--Matveev lower bound for
\[
 q\log M-p\log2=(qx-p)\log2
\]
is polynomial in $q$ (with a large exponent depending on the fixed algebraic numbers).
This is compatible with the $O(1/q)$ upper bound from continued fractions.  Synchronizing
the analogous base-three expression recreates the original quadratic determinant rather
than producing a contradiction.

\subsection{Confusing formal and numerical nonvanishing}

The tensor $q_{a,b}$ is formally nonzero for a nonintegral candidate.  It does not follow
that its specialization $e_p\mapsto\log p$ is nonzero.  Proving that implication on the
positive two-star cone is exactly \Cref{thm:depth-two-equivalence}; on all tensors it would
be a much stronger algebraic-independence statement.

\subsection{Symmetrization is a reduction, not a proof}

The equality
\[
 (\log\beta)^2=(\log\alpha)(\log\gamma)
\]
for algebraic $\alpha,\beta,\gamma$ is not ruled out by Baker when the three logarithms are
independent.  The symmetric normal form exposes a clean obstruction but does not make it
linear.

\section{Minimal missing statements}
\label{sec:minimal-missing}

\subsection{The exact two-star separation principle}

\begin{conjecture}[Positive two-star separation]
\label{conj:two-star-separation}
For nonnegative integral exponent vectors $a,b$,
\[
 \pi_2(e_2b-e_3a)=0
 \quad\Longrightarrow\quad
 e_2b-e_3a=0.
\]
\end{conjecture}

By \Cref{thm:depth-two-equivalence}, this is not merely sufficient; it is equivalent to the
Alaoglu--Erd\H{o}s conjecture.  Its advantage is diagnostic: the positivity and two-star
support are explicit, and the only allowed formal zero is visible coefficient by
coefficient.

\subsection{Rank-separated versions}

Let $\mathcal C_{2,3}^{(r)}$ be the subset of the two-star cone whose underlying exponent
span has dimension $r$.

\begin{conjecture}[Rank-three two-star separation]
\label{conj:rank3}
$\ker\pi_2\cap\mathcal C_{2,3}^{(3)}=\{0\}$.
\end{conjecture}

\begin{conjecture}[Rank-four two-star separation]
\label{conj:rank4}
$\ker\pi_2\cap\mathcal C_{2,3}^{(4)}=\{0\}$.
\end{conjecture}

\begin{theorem}[Exact branch decomposition]
\label{thm:branch-decomposition}
The Alaoglu--Erd\H{o}s conjecture is equivalent to the conjunction of
\Cref{conj:rank3,conj:rank4}.  The rank-three conjecture is equivalent to excluding the
nondegenerate ternary conics \eqref{eq:ternary-q} at the prime-log point; the rank-four
conjecture excludes full-rank hyperbolic pullbacks.
\end{theorem}

\begin{proof}
Theorem~\ref{thm:counterexample-rank} proves that every counterexample lies in exactly one of
the two branches.  Conversely, any nonzero element in either kernel intersection constructs
a counterexample exactly as in the proof of \Cref{thm:depth-two-equivalence}.  The conic
description is \Cref{sec:rank-three-conic}, and the rank-four description follows from
\Cref{thm:hyperbolic-classification}.
\end{proof}

\subsection{A Baker-completion formulation}

\begin{conjecture}[Two-star rank drop]
\label{conj:rank-drop}
If $q=e_2b-e_3a\in\mathcal C_{2,3}$ satisfies $\pi_2(q)=0$, then
$\rank(q)\le2$.
\end{conjecture}

\begin{proposition}
\label{prop:rank-drop-implies}
Conjecture~\ref{conj:rank-drop} implies the Alaoglu--Erd\H{o}s conjecture.
\end{proposition}

\begin{proof}
If $q$ is a nonzero vanishing two-star tensor, the conjectured rank drop and
\Cref{thm:baker-rank-two} contradict one another.  Thus every vanishing two-star tensor is
zero, and \Cref{thm:depth-two-equivalence} applies.
\end{proof}

This formulation isolates the only way Baker could finish the proof: one must force the
quadratic relation into the factorable range.  Theorem~\ref{thm:counterexample-rank} shows
that a counterexample would do exactly the opposite, retaining rank three or four.

\subsection{A symmetric orbit formulation for rank three}

\begin{conjecture}[No algebraic $\theta$-orbit of length three]
\label{conj:theta-orbit}
Let $\theta=\log_2 3$.  There is no algebraic $\alpha>1$ for which
$\alpha,\alpha^\theta,\alpha^{\theta^2}$ are all algebraic and multiplicatively independent.
\end{conjecture}

\begin{corollary}
Conjecture~\ref{conj:theta-orbit} rules out the multiplicative-rank-three branch of the
Alaoglu--Erd\H{o}s conjecture.
\end{corollary}

\begin{proof}
This is the contrapositive of \Cref{thm:gp-normal-form}.
\end{proof}

This is not asserted to settle rank four.  It is a sharply targeted three-number problem
that captures the entire dependent branch.

\section{Research directions suggested by the reduction}
\label{sec:research-directions}

\subsection{Attack rank three first}

The rank-three branch has several extra structures absent in rank four:
\begin{itemize}
\item one common external prime direction $R$;
\item a split nondegenerate ternary conic with a rational point;
\item the explicit determinant \eqref{eq:H-determinant};
\item the rational-function relation \eqref{eq:rho-rational-function};
\item the symmetric orbit $\alpha,\alpha^\theta,\alpha^{\theta^2}$.
\end{itemize}
A promising theorem would exploit these simultaneously rather than treating an arbitrary
Four Exponentials matrix.  In particular, any auxiliary-polynomial construction can be
specialized to a conic with known rational parametrization and to nonnegative integral
support data.

\subsection{A tensor-square version of the one-period theory}

The modern one-period results explain all linear relations among the first-level periods.
The present reduction asks for a highly restricted statement in the tensor square of the
Tate logarithm part.  One possible program is:
\begin{enumerate}[label=\textup{(\arabic*)}]
\item identify the motivic tensor represented by $e_2b-e_3a$;
\item classify functorial or shuffle-type relations that can meet the two-star cone;
\item prove that every such motivic relation is already the formal zero of
\Cref{lem:formal-zero};
\item use an appropriate period injectivity theorem for this restricted tensor class.
\end{enumerate}
The missing step is not expected to follow from the existing linear classification alone.

\subsection{Support-adapted matrix-coefficient constructions}

The general Matrix Coefficient Conjecture permits arbitrary logarithm matrices.  Here every
entry comes from two fixed prime rows and two nonnegative integer exponent columns.  In the
prime basis, the determinant has only two star centers.  An auxiliary polynomial or
fewnomial argument that preserves this support may need substantially fewer conditions than
the general matrix problem.  The target should be a vanishing rational matrix coefficient,
not merely another small determinant estimate.

\subsection{Formal verification targets}

The new linear algebra is well suited to formalization.  A compact verified development
could include:
\begin{itemize}
\item the map $\pi_2$ and formal zero classification;
\item rank-one/rank-two factorization over an algebraic closure;
\item the hyperbolic restriction calculation and tangency invariant;
\item the determinant formula \eqref{eq:H-determinant};
\item rational symmetrization by $Q^{-T}=C^{-1}J$.
\end{itemize}
Formalization would certify the reductions, not the missing transcendence statement.

\section{Conclusions}
\label{sec:conclusion}

The Alaoglu--Erd\H{o}s statement is elementary to formulate and extraordinarily rigid, but
its exact difficulty is modern transcendence theory rather than elementary divisibility.
The investigation in this report yields the following clean picture.

First, any counterexample is a positive transcendental exponent and creates a singular
$2\times2$ matrix of algebraic logarithms with no vanishing rational matrix coefficient.
That is exactly the size-two Matrix Coefficient/Four Exponentials obstruction.

Second, in the formal prime-exponent space the determinant is the sparse quadratic tensor
$q_{a,b}=e_2b-e_3a$.  The conjecture is exactly injectivity of the degree-two period map on
the positive two-star cone.

Third, Baker's theorem goes farther than a purely linear audit might suggest: it excludes
every vanishing two-star tensor of rank one or two.  Therefore a counterexample must carry a
genuine rank-three or rank-four quadratic relation.  The hyperbolic pullback theorem
classifies the ranks exactly.

Fourth, the rank-three branch admits a new symmetric normal form.  It would force an
algebraic three-term orbit
$\alpha,\alpha^{\log_2 3},\alpha^{(\log_2 3)^2}$ with multiplicatively independent terms.
This isolates a compact special obstruction, but current theorems do not exclude it.

The remaining mathematical task is therefore precise: prove nonvanishing for the rank-three
and rank-four positive two-star tensors, or prove the corresponding vanishing matrix
coefficient.  Doing so would prove the conjecture.  This report establishes the reductions
and the low-rank theorem, but it does not pretend that the final depth-two period statement
has already been proved.

\appendix

\section{Coordinate proof of the hyperbolic restriction criterion}
\label{app:hyperbolic}

For reference, we expand the rank-three calculation in coordinates.  Let
\[
 G=\begin{pmatrix}
 0&1&0&0\\1&0&0&0\\0&0&0&-1\\0&0&-1&0
 \end{pmatrix},
 \qquad c=(c_1,c_2,c_3,c_4)^T.
\]
The hyperplane $H=c^\perp$ has dimension three.  A vector $z\in H$ belongs to the radical of
$G|_H$ exactly when $z^TGh=0$ for every $h\in H$.  With the ordinary coordinate pairing,
this means $Gz$ belongs to the annihilator of $H$, namely $\Q c$.  Since $G^{-1}=G$, such a
vector must be $z=tGc$.  It lies in $H$ exactly when
\[
 0=c^Tz=t\,c^TGc
 =2t(c_1c_2-c_3c_4).
\]
If $c_1c_2-c_3c_4\ne0$, the only radical vector is zero and the restriction has rank three.
If it vanishes, $Gc$ is a nonzero radical vector.  The radical cannot have dimension more
than one because $G$ is nondegenerate and $H$ has codimension one, so the restricted rank is
two.

\section{Derivation of the rank-three determinant}
\label{app:determinant}

Starting with
\[
 Q=CX^2+(D-A)XY-BY^2+vXZ-uYZ,
\]
twice its symmetric matrix is
\[
 H=\begin{pmatrix}
 2C&D-A&v\\D-A&-2B&-u\\v&-u&0
 \end{pmatrix}.
\]
Expanding along the last row gives
\begin{align*}
 \det H
 &=v\det\begin{pmatrix}D-A&v\\-2B&-u\end{pmatrix}
   +u\det\begin{pmatrix}2C&v\\D-A&-u\end{pmatrix}\\
 &=v\bigl(-(D-A)u+2Bv\bigr)
   +u\bigl(-2Cu-v(D-A)\bigr)\\
 &=2\bigl((A-D)uv+Bv^2-Cu^2\bigr).
\end{align*}
The point $[0:0:1]$ lies on $Q=0$.  On the affine chart $X\ne0$, setting
$\theta=Y/X$ and $\rho=Z/X$ gives
\[
 C+(D-A)\theta-B\theta^2+(v-u\theta)\rho=0,
\]
which is \eqref{eq:rho-rational-function}.

\section{One-page conditional proof under Matrix Coefficient}
\label{app:conditional-proof}

Assume the $2\times2$ Matrix Coefficient Conjecture.  Let $x\in\R$ satisfy
$M=2^x,N=3^x\in\Z$.

If $x<0$, then $0<M<1$, impossible.  If $x\in\Q$, unique factorization applied to
$2^x=M$ gives $x\in\Z$.  It remains to exclude irrational $x$.

Form
\[
 L=\begin{pmatrix}\log2&\log M\\\log3&\log N\end{pmatrix}
  =\begin{pmatrix}\log2\\\log3\end{pmatrix}(1,x).
\]
Then $\det L=0$.  By Matrix Coefficient, there are nonzero $w,z\in\Q^2$ such that
$w^TLz=0$.  But
\[
 w^TLz=(w_1\log2+w_2\log3)(z_1+xz_2).
\]
The first factor is nonzero by unique factorization of $2$ and $3$.  The second is nonzero
because $x$ is irrational.  This contradiction proves $x\in\Z$.

Every line after the invocation of Matrix Coefficient is elementary.  The invocation is the
unproved step, equivalent in size two to Four Exponentials.

\section{Nonduplication map}
\label{app:nonduplication}

The following table records the intended separation between the supplied unified report and
the present continuation.

\begin{longtable}{@{}>{\raggedright\arraybackslash}p{0.31\textwidth}>{\raggedright\arraybackslash}p{0.61\textwidth}@{}}
\toprule
Theme already extensively represented in the supplied report & Treatment here \\
\midrule
\endfirsthead
\toprule
Theme already extensively represented in the supplied report & Treatment here \\
\midrule
\endhead
Exact solution monoid, primitive generator, counting, density & Not repeated. \\
Valuation lattices, radical and smooth-output formulations & Used only for elementary prime-exponent notation; no repetition of the lattice theory. \\
Determinant divisibility and interpolation constructions & Not repeated. \\
$p$-adic, tropical, Mahler, probability, matrix, and special-function routes & Not repeated. \\
Four Exponentials and restricted Matrix Coefficient reduction & Recalled only as the audited frontier needed to place the new results. \\
Prime-log product independence as a sufficient condition & Replaced by the finer rank-sensitive Baker theorem, which proves nonvanishing for all rank-$\le2$ quadratics without assuming general product independence. \\
Multiplicative dependence/PGL branching & Reinterpreted through the hyperbolic tangency invariant and then rationally symmetrized. \\
Tower formulations & Not redeveloped; only the new symmetric consequence $\alpha,\alpha^\theta,\alpha^{\theta^2}$ arising from rank-three symmetrization is stated. \\
Formal Lean corpus & Audited conceptually; no claims of a formal proof of the conjecture. \\
\bottomrule
\caption{How the present report avoids duplicating the supplied report.}
\end{longtable}

\begin{thebibliography}{99}

\bibitem{AlaogluErdos1944}
L. Alaoglu and P. Erd\H{o}s,
\emph{On highly composite and similar numbers},
Transactions of the American Mathematical Society \textbf{56} (1944), 448--469.

\bibitem{Baker1966}
A. Baker,
\emph{Linear forms in the logarithms of algebraic numbers I, II, III},
Mathematika \textbf{13} (1966), 204--216; \textbf{14} (1967), 102--107 and 220--228.

\bibitem{Baker1975}
A. Baker,
\emph{Transcendental Number Theory},
Cambridge University Press, 1975.

\bibitem{Dasgupta2023}
S. Dasgupta,
\emph{Ranks of matrices of logarithms of algebraic numbers I},
arXiv:2303.02037, 2023.

\bibitem{DK2026}
S. Dasgupta and M. Kakde,
\emph{Ranks of matrices of logarithms of algebraic numbers II: The Matrix Coefficient Conjecture},
Advances in Mathematics \textbf{487} (2026), Article 110753;
DOI: 10.1016/j.aim.2025.110753; arXiv:2408.08178.

\bibitem{Gelfond1934}
A. O. Gel'fond,
\emph{Sur le septi\`eme probl\`eme de Hilbert},
Izvestiya Akademii Nauk SSSR \textbf{7} (1934), 623--630.

\bibitem{HW2022}
A. Huber and G. W\"ustholz,
\emph{Transcendence and Linear Relations of 1-Periods},
Cambridge Tracts in Mathematics 227, Cambridge University Press, 2022.

\bibitem{HuberKalck2026}
A. Huber and M. Kalck,
\emph{Dimension formulas for period spaces via motives and species},
arXiv:2405.21053v2 (27 March 2025; manuscript dated 11 August 2026).

\bibitem{Lang1966}
S. Lang,
\emph{Introduction to Transcendental Numbers},
Addison--Wesley, 1966.

\bibitem{Ramachandra1967}
K. Ramachandra,
\emph{Contributions to the theory of transcendental numbers I, II},
Acta Arithmetica \textbf{14} (1967/68), 65--72 and 73--88.

\bibitem{Roy1992}
D. Roy,
\emph{Matrices whose coefficients are linear forms in logarithms},
Journal of Number Theory \textbf{41} (1992), 22--47.

\bibitem{Schneider1957}
T. Schneider,
\emph{Einf\"uhrung in die transzendenten Zahlen},
Springer, 1957.

\bibitem{SertozOuaknineWorrell2025}
E. C. Sert\"oz, J. Ouaknine, and J. Worrell,
\emph{Computing transcendence and linear relations of 1-periods},
arXiv:2505.20397, 2025.

\bibitem{Waldschmidt2000}
M. Waldschmidt,
\emph{Diophantine Approximation on Linear Algebraic Groups},
Grundlehren der mathematischen Wissenschaften 326, Springer, 2000.

\bibitem{Waldschmidt2022}
M. Waldschmidt,
\emph{The Four Exponentials Problem and Schanuel's Conjecture},
in \emph{Mathematics Going Forward}, Lecture Notes in Mathematics 2313,
Springer, 2022, 579--592.

\end{thebibliography}

\end{document}
